% GENERATED FILE. Edit canonical lessons, then run npm run generate:books.
% canonical-source-hash: fnv1a64:b4d88015d2e98fcb
% canonical-lessons: ML-C414-tatvam, ML-C414-dharana, ML-C414-siddhantam, ML-C414-uravitam, ML-C414-satamanam

\chapter{Ideas and Sources}
\label{ch:ml-ideas-and-sources}
\hlchaptermodality{414}

\begin{chapteropening}
\textbf{By the end of this chapter:} \emph{I can say a principle, an understanding, a theory, a source and per cent.}

Complete the last lesson of chapter 414: I can say a principle, an understanding, a theory, a source and per cent.
\end{chapteropening}

\section[\texorpdfstring{\ml{തത്വം} (tatvaṁ)}{തത്വം (tatvaṁ)}]{\ml{തത്വം} (tatvaṁ) — a principle}

\label{lesson:ML-C414-tatvam}



\begin{quote}
Before the new one: say the Malayalam for to urge, then the Malayalam for to express.
\end{quote}

\subsection*{You'll want to know: \ml{തത്വം}}
\textbf{\ml{തത്വം}} — \emph{tatvaṁ} — \textquotedblleft{}a principle\textquotedblright{}.

\begin{quote}
\ml{സത്യം} \ml{പറയുക} \ml{എന്നതാണ്} \ml{അവന്റെ} \ml{തത്വം}.
\end{quote}

\begin{quote}
\emph{satyaṁ paṟayuka ennatāṇŭ avanṟe tatvaṁ.}
\end{quote}

\begin{quote}
Telling the truth is his principle.
\end{quote}

Also spelled \ml{തത്ത്വം} in older books.

\begin{cousinweb}
\textbf{\ml{തത്വം}} (\emph{tatvaṁ}) is borrowed from Sanskrit \emph{tattva}, \textquotedblleft{}the real state of a thing, truth\textquotedblright{} (CDIAL 5642), from \emph{tad}, \textquotedblleft{}that\textquotedblright{}: thatness, what a thing is.
\end{cousinweb}

\begin{grammarlens}[title={What you've built}]
One more word for giving reasons and opinions.
\end{grammarlens}

\subsection*{Your turn}
\begin{itemize}
\raggedright
  \item \emph{Say it:} \emph{tatvaṁ}
  \item \emph{Say it:} \emph{tatvaṁ}, once more
  \item \emph{Say it:} \emph{satyaṁ paṟayuka ennatāṇŭ avanṟe tatvaṁ.} — Telling the truth is his principle.
  \item \emph{Recall:} say the Malayalam for to urge, then the Malayalam for to express, then say \emph{tatvaṁ} again
\end{itemize}

\begin{tcolorbox}[breakable,colback=teal!4,colframe=teal!35!black,title={Before you move on}]
What does \ml{തത്വം} mean? (\textquotedblleft{}A principle\textquotedblright{}.) Say it once more.
\end{tcolorbox}
\section[\texorpdfstring{\ml{ധാരണ} (dhāraṇa)}{ധാരണ (dhāraṇa)}]{\ml{ധാരണ} (dhāraṇa) — an understanding; a notion}

\label{lesson:ML-C414-dharana}



\begin{quote}
Before the new one: say the Malayalam for to express, then the Malayalam for a principle.
\end{quote}

\subsection*{You'll want to know: \ml{ധാരണ}}
\textbf{\ml{ധാരണ}} — \emph{dhāraṇa} — \textquotedblleft{}an understanding; a notion\textquotedblright{}.

\begin{quote}
\ml{ആ} \ml{വിഷയത്തെക്കുറിച്ച്} \ml{എനിക്ക്} \ml{നല്ല} \ml{ധാരണയുണ്ട്}.
\end{quote}

\begin{quote}
\emph{ā viṣayattekkuṟiccŭ enikkŭ nalla dhāraṇayuṇṭŭ.}
\end{quote}

\begin{quote}
I have a good understanding of that subject.
\end{quote}

It is also an agreement between two sides: \ml{അവർ} \ml{തമ്മിൽ} \ml{ധാരണയായി}, they reached an understanding.

\begin{cousinweb}
\textbf{\ml{ധാരണ}} (\emph{dhāraṇa}) is borrowed from Sanskrit \emph{dhāraṇā}, \textquotedblleft{}holding, bearing; keeping in the mind\textquotedblright{} (Monier-Williams). An understanding is what the mind holds.
\end{cousinweb}

\begin{grammarlens}[title={What you've built}]
One more word for giving reasons and opinions.
\end{grammarlens}

\subsection*{Your turn}
\begin{itemize}
\raggedright
  \item \emph{Say it:} \emph{dhāraṇa}
  \item \emph{Say it:} \emph{dhāraṇa}, once more
  \item \emph{Say it:} \emph{ā viṣayattekkuṟiccŭ enikkŭ nalla dhāraṇayuṇṭŭ.} — I have a good understanding of that subject.
  \item \emph{Recall:} say the Malayalam for to express, then the Malayalam for a principle, then say \emph{dhāraṇa} again
\end{itemize}

\begin{tcolorbox}[breakable,colback=teal!4,colframe=teal!35!black,title={Before you move on}]
What does \ml{ധാരണ} mean? (\textquotedblleft{}An understanding; a notion\textquotedblright{}.) Say it once more.
\end{tcolorbox}
\section[\texorpdfstring{\ml{സിദ്ധാന്തം} (siddhāntaṁ)}{സിദ്ധാന്തം (siddhāntaṁ)}]{\ml{സിദ്ധാന്തം} (siddhāntaṁ) — a theory}

\label{lesson:ML-C414-siddhantam}



\begin{quote}
Before the new one: say the Malayalam for a principle, then the Malayalam for an understanding.
\end{quote}

\subsection*{You'll want to know: \ml{സിദ്ധാന്തം}}
\textbf{\ml{സിദ്ധാന്തം}} — \emph{siddhāntaṁ} — \textquotedblleft{}a theory\textquotedblright{}.

\begin{quote}
\ml{പുതിയ} \ml{സിദ്ധാന്തം} \ml{പലരും} \ml{അംഗീകരിച്ചില്ല}.
\end{quote}

\begin{quote}
\emph{putiya siddhāntaṁ palaruṁ aṁgīkariccilla.}
\end{quote}

\begin{quote}
Many people did not accept the new theory.
\end{quote}

Plural: \ml{സിദ്ധാന്തങ്ങൾ}, \emph{siddhāntaṅṅaḷ}.

Plural: \ml{സിദ്ധാന്തങ്ങൾ}.

\begin{cousinweb}
\textbf{\ml{സിദ്ധാന്തം}} (\emph{siddhāntaṁ}) is borrowed from Sanskrit \emph{siddhānta}, \textquotedblleft{}an established end, a doctrine\textquotedblright{} (CDIAL 13404): \emph{siddha}, \textquotedblleft{}settled\textquotedblright{}, and \emph{anta}, \textquotedblleft{}end\textquotedblright{}.
\end{cousinweb}

\begin{grammarlens}[title={What you've built}]
One more word for giving reasons and opinions.
\end{grammarlens}

\subsection*{Your turn}
\begin{itemize}
\raggedright
  \item \emph{Say it:} \emph{siddhāntaṁ}
  \item \emph{Say it:} \emph{siddhāntaṁ}, once more
  \item \emph{Say it:} \emph{putiya siddhāntaṁ palaruṁ aṁgīkariccilla.} — Many people did not accept the new theory.
  \item \emph{Recall:} say the Malayalam for a principle, then the Malayalam for an understanding, then say \emph{siddhāntaṁ} again
\end{itemize}

\begin{tcolorbox}[breakable,colback=teal!4,colframe=teal!35!black,title={Before you move on}]
What does \ml{സിദ്ധാന്തം} mean? (\textquotedblleft{}A theory\textquotedblright{}.) Say it once more.
\end{tcolorbox}
\section[\texorpdfstring{\ml{ഉറവിടം} (uṟaviṭaṁ)}{ഉറവിടം (uṟaviṭaṁ)}]{\ml{ഉറവിടം} (uṟaviṭaṁ) — a source}

\label{lesson:ML-C414-uravitam}



\begin{quote}
Before the new one: say the Malayalam for an understanding, then the Malayalam for a theory.
\end{quote}

\subsection*{You'll want to know: \ml{ഉറവിടം}}
\textbf{\ml{ഉറവിടം}} — \emph{uṟaviṭaṁ} — \textquotedblleft{}a source\textquotedblright{}.

\begin{quote}
\ml{വാർത്തയുടെ} \ml{ഉറവിടം} \ml{ആരും} \ml{പറഞ്ഞില്ല}.
\end{quote}

\begin{quote}
\emph{vārttayuṭe uṟaviṭaṁ āruṁ paṟaññilla.}
\end{quote}

\begin{quote}
Nobody said what the source of the news was.
\end{quote}

It is first the place a spring rises, then any source of news or money.

\begin{cousinweb}
\textbf{\ml{ഉറവിടം}} (\emph{uṟaviṭaṁ}) is a native Malayalam compound: \ml{ഉറവ}, \textquotedblleft{}a spring of water\textquotedblright{}, from \ml{ഊറുക}, \textquotedblleft{}to spring, to ooze\textquotedblright{} (DEDR 761), and \ml{ഇടം}, \textquotedblleft{}a place\textquotedblright{} (DEDR 434). A source is where things spring from.
\end{cousinweb}

\begin{grammarlens}[title={What you've built}]
One more word for giving reasons and opinions.
\end{grammarlens}

\subsection*{Your turn}
\begin{itemize}
\raggedright
  \item \emph{Say it:} \emph{uṟaviṭaṁ}
  \item \emph{Say it:} \emph{uṟaviṭaṁ}, once more
  \item \emph{Say it:} \emph{vārttayuṭe uṟaviṭaṁ āruṁ paṟaññilla.} — Nobody said what the source of the news was.
  \item \emph{Recall:} say the Malayalam for an understanding, then the Malayalam for a theory, then say \emph{uṟaviṭaṁ} again
\end{itemize}

\begin{tcolorbox}[breakable,colback=teal!4,colframe=teal!35!black,title={Before you move on}]
What does \ml{ഉറവിടം} mean? (\textquotedblleft{}A source\textquotedblright{}.) Say it once more.
\end{tcolorbox}
\section[\texorpdfstring{\ml{ശതമാനം} (śatamānaṁ)}{ശതമാനം (śatamānaṁ)}]{\ml{ശതമാനം} (śatamānaṁ) — per cent; a percentage}

\label{lesson:ML-C414-satamanam}



\begin{quote}
Before the new one: say the Malayalam for a theory, then the Malayalam for a source.
\end{quote}

\subsection*{You'll want to know: \ml{ശതമാനം}}
\textbf{\ml{ശതമാനം}} — \emph{śatamānaṁ} — \textquotedblleft{}per cent; a percentage\textquotedblright{}.

\begin{quote}
\ml{എഴുപത്} \ml{ശതമാനം} \ml{കുട്ടികളും} \ml{പരീക്ഷ} \ml{ജയിച്ചു}.
\end{quote}

\begin{quote}
\emph{eḻupatŭ śatamānaṁ kuṭṭikaḷuṁ parīkṣa jayiccu.}
\end{quote}

\begin{quote}
Seventy per cent of the children passed the exam.
\end{quote}

The number comes first: \ml{പത്ത്} \ml{ശതമാനം}, ten per cent.

\begin{cousinweb}
\textbf{\ml{ശതമാനം}} (\emph{śatamānaṁ}) is built from Sanskrit \emph{śata}, \textquotedblleft{}a hundred\textquotedblright{} (CDIAL 12278), and \emph{māna}, \textquotedblleft{}a measure\textquotedblright{} (CDIAL 10041): a measure out of a hundred.
\end{cousinweb}

\begin{grammarlens}[title={What you've built}]
One more word for giving reasons and opinions.
\end{grammarlens}

\subsection*{Your turn}
\begin{itemize}
\raggedright
  \item \emph{Say it:} \emph{śatamānaṁ}
  \item \emph{Say it:} \emph{śatamānaṁ}, once more
  \item \emph{Say it:} \emph{eḻupatŭ śatamānaṁ kuṭṭikaḷuṁ parīkṣa jayiccu.} — Seventy per cent of the children passed the exam.
  \item \emph{Recall:} say the Malayalam for a theory, then the Malayalam for a source, then say \emph{śatamānaṁ} again
\end{itemize}

\begin{tcolorbox}[breakable,colback=teal!4,colframe=teal!35!black,title={Before you move on}]
What does \ml{ശതമാനം} mean? (\textquotedblleft{}Per cent; a percentage\textquotedblright{}.) Say it once more.
\end{tcolorbox}
